\documentclass[11pt,letterpaper]{article}

\usepackage[margin=1in]{geometry}
\usepackage[T1]{fontenc}
\usepackage[utf8]{inputenc}
\usepackage{lmodern}
\usepackage{microtype}
\usepackage[hidelinks]{hyperref}
\usepackage{enumitem}

\setlength{\parindent}{0pt}
\setlength{\parskip}{6pt}
\setlist[itemize]{leftmargin=1.4em, itemsep=3pt, topsep=3pt}
\setlist[enumerate]{leftmargin=1.6em, itemsep=3pt, topsep=3pt}

\newcommand{\pkg}[1]{\textsf{#1}}
\newcommand{\proglang}[1]{\textsf{#1}}
\newcommand{\code}[1]{\texttt{#1}}
\newcommand{\commentheading}[1]{\vspace{8pt}\noindent\textbf{Comment #1}\par}

\title{\textbf{Response to the Editorial Return}\\
\pkg{exdqlm}: An R Package for Estimation and Analysis of Flexible Dynamic Quantile Linear Models}
\author{Antonio De Leon, Raquel Barata, Raquel Prado, and Bruno Sans\'o}
\date{Prepared September 2026}

\begin{document}
\maketitle

Dear Editors,

Thank you for the careful second prescreening review. We understand the return
as a reproducibility-interface and output-provenance problem: the previous
archive still made it too hard to run the article computations in batch mode and
compare the resulting output directly with the manuscript. We have restructured
the replication materials accordingly and are resubmitting as a new submission
with a point-by-point response.

The submitted files are:
\begin{itemize}
\item \code{exdqlm-jss.pdf}: manuscript PDF.
\item \code{exdqlm\_1.1.1.tar.gz}: software source.
\item \code{exdqlm-jss-replication.tar.gz}: replication archive.
\item \code{response-to-editor.pdf}: this response.
\end{itemize}

The package source is version 1.1.1 of \pkg{exdqlm}. This is a narrow
reproducibility patch for the manuscript rerun: it does not change the exported
API or manuscript conclusions.

\section*{Summary of the Revision}

\begin{itemize}
\item We replaced the previous replication interface with one flat
  reviewer-facing \code{code.R}. The script follows the manuscript examples in
  order and uses the same fitted objects, plotting commands, and table
  calculations. It does not call \code{source()}, read or write saved fit
  objects, inspect Git metadata, or require local paths.
\item The primary command is the JSS-style batch command
  \code{R CMD BATCH --vanilla code.R code.Rout}, with native thread variables
  set before \proglang{R} starts. The run produces \code{code.Rout} and
  \code{Rplots.pdf} at the archive root.
\item The full \code{code.Rout} prints the fitted object \code{M95},
  \code{summary(M95)}, \code{MTF\$median.kt}, Tables 7--10, and
  \code{sessionInfo()}. This makes the numerical output visible in the batch
  output instead of only in separate files.
\item The graphics stream produced by \code{code.R} contains the manuscript
  figures in manuscript order. The same plotting expressions also write the PNG
  files used by the manuscript.
\item The manuscript now includes a computational-details paragraph describing
  the reference \proglang{R} version, \pkg{exdqlm} version, RNG settings,
  backend/thread settings, and the role of the full replication script.
\item The README and reproducibility notes now describe the public replication
  command and expected outputs. Development-only materials are not included as
  reviewer-facing entry points.
\end{itemize}

\section*{Responses to the Specific Comments}

\commentheading{1. Single script and batch-mode output}

The revised archive centers on one \proglang{R} script, \code{code.R}. It is a
flat, commented file ordered by the manuscript examples. Running
\code{R CMD BATCH --vanilla code.R code.Rout} performs the full manuscript
replication and creates the batch output file requested in the editorial
comment.

We removed the ambiguous \code{examples.R} entry point from the JSS-facing
archive. The submitted archive now gives reviewers one obvious script to run
and one batch output file to inspect.

\commentheading{2. Exact reproducibility and computing environment}

The public script sets
\[
\code{RNGversion("4.6.0")}
\]
and
\[
\code{RNGkind("Mersenne-Twister", "Inversion", "Rejection")}.
\]
It also stops immediately unless the loaded package is \pkg{exdqlm} version
1.1.1.

While investigating the differences reported by the editors, we also tightened
the package's stochastic reproducibility controls. Version 1.1.1 is the public
package version used for the resubmission, and the full JSS replication was
rerun with this version.

The manuscript and README now state the reference environment and the thread
variables used before launching \proglang{R}. Runtime values are still described
as platform-dependent elapsed fitting times.

\commentheading{3. More output in the manuscript and batch output}

We added visible fitted-object output to the replication script and manuscript
discussion. The full batch output now includes
\[
\code{M95}
\]
and
\[
\code{summary(M95)}.
\]
The script also prints Tables 7--10 and \code{MTF\$median.kt} directly to
\code{code.Rout}. This allows a reviewer to inspect the generated batch output
without searching through intermediate files.

\commentheading{4. Figure 1 and Figure 9 differences}

Both figures are now regenerated by \code{code.R} from the same plotting
commands used to write the manuscript PNG files. The revised archive makes the
figure provenance direct and rerunnable from the submitted script.

\commentheading{5. \code{MTF\$median.kt}}

The manuscript value is now taken from the full \code{code.R} reference run.
The scalar is printed explicitly in \code{code.Rout}, so it can be compared
directly with the manuscript.

\commentheading{6. Missing Figure 3 top panel and Figure 5}

\code{code.R} now regenerates the full graphics stream in manuscript order.
The stream includes the complete Sunspots figure with its top panel and the Big
Tree data/covariate figure. The corresponding PNG files are written under
\code{figures/}.

\commentheading{7. Missing Tables 7 and 10}

Tables 7 and 10 are now generated by the full script and printed to
\code{code.Rout}. The same table objects are also saved as CSV files under
\code{tables/}.

\commentheading{8. \code{tab.ex3} and \code{tab.ex3.fc}}

The Big Tree fitted-diagnostic table and held-out forecast table are now printed
to \code{code.Rout} as Tables 8 and 9. They are also saved as CSV files in the
generated table directory.

\commentheading{9. Fitted-object output}

The full script prints the \code{M95} object and \code{summary(M95)}. The
manuscript now guides the reader through representative fitted-object output,
so the standard \code{print()} and \code{summary()} methods are visible in the
article and in the batch output.

\section*{Validation}

Before resubmission we ran the following checks with \proglang{R} 4.6.0:
\begin{itemize}
\item targeted package repeatability tests for compiled samplers, dynamic MCMC,
  and static MCMC;
\item the package test suite and CRAN submission checks for
  \code{exdqlm\_1.1.1.tar.gz};
\item \code{R CMD BATCH --vanilla code.R code.Rout};
\item manuscript and response compilation; and
\item archive extraction, checksum verification, script parsing, package-version
  verification, and manuscript-source compilation from a directory without Git
  metadata.
\end{itemize}

We believe the revised materials now provide a simpler and more standard JSS
replication interface, with the numerical and graphical outputs visible in the
batch files requested by the editors.

Sincerely,

Antonio De Leon, Raquel Barata, Raquel Prado, and Bruno Sans\'o

\end{document}
